\documentclass[11pt,a4paper,fontset=none]{ctexart}

\usepackage[margin=2.2cm]{geometry}
\usepackage{amsmath,amssymb,bm}
\usepackage{booktabs}
\usepackage{longtable}
\usepackage{array}
\usepackage{xcolor}
\usepackage{hyperref}
\usepackage{enumitem}

\setCJKmainfont{Songti SC}
\setCJKsansfont{Songti SC}
\setCJKmonofont{Songti SC}

\hypersetup{
  colorlinks=true,
  linkcolor=blue!50!black,
  citecolor=blue!50!black,
  urlcolor=blue!50!black
}

\setlist[itemize]{leftmargin=1.7em}
\setlist[enumerate]{leftmargin=1.9em}

\title{\bfseries 超表面深度成像项目公式推导整理\\
\large 从 RGB-D 数据、点源 PSF、整图成像到 Fisher 信息优化}
\author{项目推导笔记}
\date{2026年6月}

\begin{document}
\maketitle
\tableofcontents
\newpage

\section{总览}

本项目的核心链路可以写成：
\begin{equation}
  I_{\mathrm{rgb}},D,K
  \longrightarrow
  \text{场景点源集合}
  \longrightarrow
  h(z,\alpha,\lambda;\phi,d)
  \longrightarrow
  Y_{\mathrm{CMOS}}
  \longrightarrow
  \hat D .
\end{equation}

其中 $I_{\mathrm{rgb}}$ 是普通 RGB 图像，$D$ 是对应深度图，$K$ 是相机内参，$h$ 是一个点经过超表面后落到 CMOS 上形成的点扩散函数（PSF），$Y_{\mathrm{CMOS}}$ 是超表面相机观测图像，$\hat D$ 是恢复出的深度图。

第一阶段建议采用单波长窄带模型：
\begin{equation}
  \lambda=\lambda_0,
\end{equation}
即实际系统在超表面前放置窄带滤光片，仿真中只计算 $\lambda_0$ 附近的光传播。这可以显著降低 RGB 到连续光谱不可逆的问题。

\section{符号定义}

\begin{longtable}{p{3.4cm}p{10.4cm}}
\toprule
符号 & 含义 \\
\midrule
\endhead
$I_{\mathrm{rgb}}(x,y)$ & 输入 RGB 图像，像素坐标为 $(x,y)$。 \\
$I_c(x,y)$ & RGB 第 $c$ 个通道，$c\in\{R,G,B\}$。 \\
$I_c^{\mathrm{lin}}(x,y)$ & 线性光强域的 RGB 通道。 \\
$D(x,y)$ & 深度图。可表示 optical-axis depth 或 range depth，必须按数据集定义确认。 \\
$K$ & 相机内参矩阵。 \\
$f_x,f_y$ & 以像素为单位的水平、垂直焦距。 \\
$c_x,c_y$ & 主点坐标。 \\
$\alpha_x,\alpha_y$ & 像素对应光线的水平、垂直入射角。 \\
$\lambda$ & 波长。第一阶段固定为 $\lambda_0$。 \\
$k$ & 波数，$k=2\pi/\lambda$。 \\
$(x_m,y_m)$ & 超表面平面坐标。 \\
$(x_i,y_i)$ 或 $(u,v)$ & CMOS 平面坐标。 \\
$d$ & 超表面到 CMOS 的传播距离。 \\
$P(x_m,y_m)$ & 超表面孔径函数。 \\
$t_\phi(x_m,y_m,\lambda)$ & 超表面复透射函数。 \\
$\phi(x_m,y_m,\lambda)$ & 超表面相位分布。 \\
$U_{\mathrm{in}}$ & 入射到超表面的复振幅。 \\
$U_0$ & 经过超表面调制后的复振幅。 \\
$U_i$ & 到达 CMOS 平面的复振幅。 \\
$h$ & 点扩散函数 PSF，$h=|U_i|^2$。 \\
$Y(u,v)$ & CMOS 上理想光强或期望响应。 \\
$\tilde Y(u,v)$ & 加噪声和量化后的 CMOS 数字图像。 \\
$\bm\mu$ & 观测图像的均值向量。 \\
$\bm\Sigma$ & 噪声协方差矩阵。 \\
$\bm F$ & Fisher information matrix。 \\
$\mathcal I_z$ & 对深度 $z$ 的 Fisher information。 \\
\bottomrule
\end{longtable}

\section{相机内参：像素如何变成光线方向}

\subsection{针孔相机模型}

相机内参矩阵为：
\begin{equation}
  K=
  \begin{bmatrix}
  f_x&0&c_x\\
  0&f_y&c_y\\
  0&0&1
  \end{bmatrix}.
\end{equation}

三维点 $(X,Y,Z)$ 投影到像素 $(x,y)$：
\begin{equation}
  x=f_x\frac{X}{Z}+c_x,\qquad
  y=f_y\frac{Y}{Z}+c_y.
\end{equation}

反过来，像素坐标对应的相机光线方向为：
\begin{equation}
  \bm r(x,y)
  =
  K^{-1}
  \begin{bmatrix}
  x\\y\\1
  \end{bmatrix}
  =
  \begin{bmatrix}
  \frac{x-c_x}{f_x}\\
  \frac{y-c_y}{f_y}\\
  1
  \end{bmatrix}.
\end{equation}

单位方向向量为：
\begin{equation}
  \hat{\bm r}(x,y)
  =
  \frac{\bm r(x,y)}{\|\bm r(x,y)\|}.
\end{equation}

对应视场角为：
\begin{equation}
  \alpha_x(x,y)=
  \arctan\left(\frac{x-c_x}{f_x}\right),
  \qquad
  \alpha_y(x,y)=
  \arctan\left(\frac{y-c_y}{f_y}\right).
\end{equation}

这一步的意义是：
\begin{quote}
将图像像素坐标转换成进入超表面的物理入射方向。PSF 在大视场下依赖 $\alpha_x,\alpha_y$，因此这一步不能省略。
\end{quote}

\subsection{深度图如何变成三维点}

若深度图 $D(x,y)$ 是 optical-axis depth，即相机坐标的 $Z$，则：
\begin{equation}
  \bm P(x,y)
  =
  D(x,y)
  \begin{bmatrix}
  \frac{x-c_x}{f_x}\\
  \frac{y-c_y}{f_y}\\
  1
  \end{bmatrix}.
\end{equation}
即：
\begin{equation}
  X=D(x,y)\frac{x-c_x}{f_x},\quad
  Y=D(x,y)\frac{y-c_y}{f_y},\quad
  Z=D(x,y).
\end{equation}

若深度图是 range depth，即从相机中心到点的欧氏距离 $D_r(x,y)$，则：
\begin{equation}
  \bm P(x,y)=D_r(x,y)\hat{\bm r}(x,y).
\end{equation}

两者关系为：
\begin{equation}
  D_r=
  D_z
  \sqrt{
  \left(\frac{x-c_x}{f_x}\right)^2+
  \left(\frac{y-c_y}{f_y}\right)^2+
  1
  },
  \qquad
  D_z=D_r\hat r_z .
\end{equation}

\section{RGB 图像如何近似为超表面前的入射强度}

\subsection{sRGB 到线性光强}

普通图像通常是非线性的 sRGB。若原始像素为 $I_c^{8bit}(x,y)\in[0,255]$，先归一化：
\begin{equation}
  I_c^s(x,y)=\frac{I_c^{8bit}(x,y)}{255}.
\end{equation}

简化的 inverse gamma 近似为：
\begin{equation}
  I_c^{\mathrm{lin}}(x,y)=\left(I_c^s(x,y)\right)^{2.2}.
\end{equation}

更标准的 sRGB inverse gamma 为：
\begin{equation}
I_c^{\mathrm{lin}}=
\begin{cases}
\frac{I_c^s}{12.92}, & I_c^s\le 0.04045,\\
\left(\frac{I_c^s+0.055}{1.055}\right)^{2.4}, & I_c^s>0.04045.
\end{cases}
\end{equation}

\subsection{为什么 RGB 不能唯一恢复连续光谱}

真实光谱强度可写为 $L(x,y,\lambda)$。普通 RGB 通道本质上是连续光谱经相机光谱响应积分后的结果：
\begin{equation}
  I_R(x,y)=\int L(x,y,\lambda)Q_R(\lambda)d\lambda,
\end{equation}
\begin{equation}
  I_G(x,y)=\int L(x,y,\lambda)Q_G(\lambda)d\lambda,
\end{equation}
\begin{equation}
  I_B(x,y)=\int L(x,y,\lambda)Q_B(\lambda)d\lambda.
\end{equation}

已知三个积分值 $I_R,I_G,I_B$，无法唯一恢复连续函数 $L(\lambda)$。因此 RGB-D 数据不能完整变成真实自然光场，只能做建模近似。

\subsection{单波长窄带近似}

若系统在超表面前加入中心波长 $\lambda_0$ 的窄带滤光片，则可近似：
\begin{equation}
  L(x,y,\lambda)\approx S(x,y;\lambda_0)\delta(\lambda-\lambda_0).
\end{equation}

若 $\lambda_0$ 位于绿光附近，例如 $\lambda_0\approx 550$ nm，则第一版可用绿色通道近似该波长强度：
\begin{equation}
  S(x,y;\lambda_0)\approx \beta I_G^{\mathrm{lin}}(x,y).
\end{equation}
其中 $\beta$ 表示滤光片和系统透过率导致的光通量缩放。理论验证阶段可先令 $\beta=1$。

\subsection{RGB 三代表波长近似}

若不做单波长近似，可设：
\begin{equation}
  \lambda_R\approx 620\,\mathrm{nm},\quad
  \lambda_G\approx 550\,\mathrm{nm},\quad
  \lambda_B\approx 460\,\mathrm{nm}.
\end{equation}

用三条谱线近似 RGB：
\begin{equation}
  L(x,y,\lambda)
  \approx
  I_R^{\mathrm{lin}}(x,y)\delta(\lambda-\lambda_R)
  +I_G^{\mathrm{lin}}(x,y)\delta(\lambda-\lambda_G)
  +I_B^{\mathrm{lin}}(x,y)\delta(\lambda-\lambda_B).
\end{equation}

第一阶段建议优先使用单波长窄带模型，因为其物理和实验对应最清晰。

\section{一个点源经过超表面形成 PSF 的推导}

\subsection{系统几何}

设超表面位于 $z=0$ 平面，坐标为 $(x_m,y_m)$。CMOS 位于 $z=d$ 平面，坐标为 $(x_i,y_i)$。点源位于超表面前方，深度为 $z_o$，视场角为 $\alpha_x,\alpha_y$，则点源坐标可近似写为：
\begin{equation}
  X_o=z_o\tan\alpha_x,\qquad
  Y_o=z_o\tan\alpha_y,\qquad
  Z_o=-z_o.
\end{equation}

点源到超表面上一点的距离为：
\begin{equation}
  R_o(x_m,y_m)
  =
  \sqrt{
  (x_m-X_o)^2+
  (y_m-Y_o)^2+
  z_o^2
  }.
\end{equation}
代入角度参数：
\begin{equation}
  R_o(x_m,y_m)
  =
  \sqrt{
  (x_m-z_o\tan\alpha_x)^2+
  (y_m-z_o\tan\alpha_y)^2+
  z_o^2
  }.
\end{equation}

\subsection{点源到达超表面的入射场}

点源在超表面处产生球面波：
\begin{equation}
  U_{\mathrm{in}}(x_m,y_m)
  =
  A_o
  \frac{
  \exp\left[i k R_o(x_m,y_m)\right]
  }{
  R_o(x_m,y_m)
  },
  \qquad k=\frac{2\pi}{\lambda}.
\end{equation}

远距离近似下，可展开为平面波角度项加球面曲率项：
\begin{equation}
  U_{\mathrm{in}}(x_m,y_m)
  \approx
  A_o
  \exp
  \left[
  ik
  \left(
  x_m\sin\alpha_x+
  y_m\sin\alpha_y+
  \frac{x_m^2+y_m^2}{2z_o}
  \right)
  \right].
\end{equation}
其中线性项对应入射角，二次项对应有限距离带来的波前曲率。长距离深度困难的根源是：
\begin{equation}
  \frac{x_m^2+y_m^2}{2z_o}
\end{equation}
随 $z_o$ 变远而迅速变弱。

\subsection{超表面复透射}

标量模型下，超表面复透射函数为：
\begin{equation}
  t_\phi(x_m,y_m,\lambda)
  =
  A_\phi(x_m,y_m,\lambda)
  \exp[i\phi(x_m,y_m,\lambda)].
\end{equation}

孔径函数为：
\begin{equation}
  P(x_m,y_m)=
  \begin{cases}
  1,&(x_m,y_m)\in\Omega,\\
  0,&\text{otherwise}.
  \end{cases}
\end{equation}

经过超表面后的场为：
\begin{equation}
  U_0(x_m,y_m)
  =
  P(x_m,y_m)
  t_\phi(x_m,y_m,\lambda)
  U_{\mathrm{in}}(x_m,y_m).
\end{equation}

\subsection{从超表面传播到 CMOS}

超表面到 CMOS 的 Fresnel 传播为：
\begin{equation}
  U_i(x_i,y_i)
  =
  \frac{\exp(ikd)}{i\lambda d}
  \iint
  U_0(x_m,y_m)
  \exp
  \left[
  i\frac{k}{2d}
  \left(
  (x_i-x_m)^2+
  (y_i-y_m)^2
  \right)
  \right]
  dx_mdy_m.
\end{equation}

代入 $U_0$：
\begin{equation}
\begin{aligned}
  U_i(x_i,y_i)
  &=
  \frac{\exp(ikd)}{i\lambda d}
  \iint
  P(x_m,y_m)t_\phi(x_m,y_m,\lambda)
  U_{\mathrm{in}}(x_m,y_m)\\
  &\qquad\times
  \exp
  \left[
  i\frac{k}{2d}
  \left(
  (x_i-x_m)^2+
  (y_i-y_m)^2
  \right)
  \right]
  dx_mdy_m.
\end{aligned}
\end{equation}

\subsection{PSF 定义}

CMOS 记录强度，故 PSF 为：
\begin{equation}
  h(x_i,y_i;z_o,\alpha_x,\alpha_y,\lambda,\phi,d)
  =
  |U_i(x_i,y_i)|^2.
\end{equation}

完整写为：
\begin{equation}
\begin{aligned}
  h(x_i,y_i;z_o,\alpha_x,\alpha_y,\lambda,\phi,d)
  &=
  \left|
  \frac{\exp(ikd)}{i\lambda d}
  \iint
  P(x_m,y_m)t_\phi(x_m,y_m,\lambda)
  \frac{A_o\exp[i k R_o(x_m,y_m)]}{R_o(x_m,y_m)}
  \right.\\
  &\qquad\left.
  \times
  \exp
  \left[
  i\frac{k}{2d}
  \left(
  (x_i-x_m)^2+
  (y_i-y_m)^2
  \right)
  \right]
  dx_mdy_m
  \right|^2 .
\end{aligned}
\end{equation}

也可用传播算子简写为：
\begin{equation}
  h=
  \left|
  \mathcal P_d
  \left\{
  P\,t_\phi\,U_{\mathrm{in}}
  \right\}
  \right|^2.
\end{equation}

\subsection{角谱法与 FFT 计算}

Fresnel 或角谱传播可写为：
\begin{equation}
  U_i
  =
  \mathcal F^{-1}
  \left[
  \mathcal F(U_0)\,
  H(f_x,f_y;d)
  \right],
\end{equation}
其中角谱传播函数为：
\begin{equation}
  H(f_x,f_y;d)
  =
  \exp
  \left[
  ikd
  \sqrt{
  1-(\lambda f_x)^2-(\lambda f_y)^2
  }
  \right].
\end{equation}

\section{超表面尺寸、周期和 CMOS 距离对 PSF 的影响}

\subsection{口径决定长距离深度敏感性}

设超表面孔径半径为 $a$。远距离下，深度引起的孔径边缘相位量级为：
\begin{equation}
  \Phi_z\sim k\frac{a^2}{2z}.
\end{equation}

两个深度 $z$ 和 $z+\Delta z$ 的波前相位差约为：
\begin{equation}
  \Delta\Phi
  \approx
  k\frac{a^2}{2}
  \left|
  \frac{1}{z}-\frac{1}{z+\Delta z}
  \right|
  \approx
  k\frac{a^2\Delta z}{2z^2}.
\end{equation}

因此：
\begin{equation}
  \Delta\Phi\propto \frac{a^2}{\lambda}\frac{\Delta z}{z^2}.
\end{equation}

结论是：口径越大，长距离深度信息越强；距离越远，深度信息按 $1/z^2$ 快速减弱。

\subsection{周期数决定相位表达能力}

若超表面周期为 $p$，孔径直径为 $D_{\mathrm{ap}}=2a$，则单元数约为：
\begin{equation}
  N_{\mathrm{cell}}
  \approx
  \left(\frac{D_{\mathrm{ap}}}{p}\right)^2.
\end{equation}

周期太少时，相位图采样过粗，无法表达复杂编码；周期太大时可能出现高阶衍射。粗略要求为：
\begin{equation}
  p<
  \frac{\lambda}{n_{\mathrm{env}}(1+\sin\alpha_{\max})}.
\end{equation}

\subsection{超表面到 CMOS 距离是系统变量}

PSF 显式依赖传播距离 $d$：
\begin{equation}
  h=
  \left|
  \mathcal P_d
  \left[
  P t_\phi U_{\mathrm{in}}
  \right]
  \right|^2.
\end{equation}

因此 $d$ 改变会直接改变 PSF。$d$ 太小，编码可能未充分展开；$d$ 太大，PSF 可能过大、光能分散、SNR 下降。理论上可将 $d$ 也纳入优化：
\begin{equation}
  \phi^\star,d^\star
  =
  \arg\max_{\phi,d}
  \mathbb E_{z,\alpha}
  \left[
  \log \mathcal I_z^{\mathrm{eff}}(z,\alpha;\phi,d)
  \right].
\end{equation}

\section{从 RGB-D 到 CMOS 整图响应}

\subsection{每个像素作为一个小点源}

单波长窄带模型下，RGB-D 图被近似为许多小点源集合：
\begin{equation}
  (x,y)
  \mapsto
  \left\{
  S(x,y;\lambda_0),\,
  D(x,y),\,
  \alpha_x(x,y),\,
  \alpha_y(x,y)
  \right\}.
\end{equation}

其中：
\begin{equation}
  S(x,y;\lambda_0)\approx I_G^{\mathrm{lin}}(x,y).
\end{equation}

\subsection{通用空间变化成像公式}

一个场景像素 $(x,y)$ 对 CMOS 像素 $(u,v)$ 的贡献为：
\begin{equation}
  dY_{x,y\rightarrow u,v}
  =
  S(x,y;\lambda_0)
  h_{\lambda_0}
  \left(
  u,v;
  x,y,
  D(x,y),
  \alpha_x(x,y),
  \alpha_y(x,y)
  \right).
\end{equation}

所有像素叠加：
\begin{equation}
\boxed{
  Y(u,v)
  =
  \sum_x\sum_y
  S(x,y;\lambda_0)
  h_{\lambda_0}
  \left(
  u,v;
  x,y,
  D(x,y),
  \alpha_x(x,y),
  \alpha_y(x,y)
  \right).
}
\end{equation}

代入绿色通道近似：
\begin{equation}
\boxed{
  Y(u,v)
  =
  \sum_x\sum_y
  I_G^{\mathrm{lin}}(x,y)
  h_{\lambda_0}
  \left(
  u,v;
  x,y,
  D(x,y),
  \alpha_x(x,y),
  \alpha_y(x,y)
  \right).
}
\end{equation}

\subsection{局部平移不变近似}

若近似认为 PSF 只依赖相对位置：
\begin{equation}
  h_{\lambda_0}(u,v;x,y,D,\alpha_x,\alpha_y)
  =
  h_{\lambda_0}(u-x,v-y;D,\alpha_x,\alpha_y),
\end{equation}
则：
\begin{equation}
  Y(u,v)
  =
  \sum_x\sum_y
  I_G^{\mathrm{lin}}(x,y)
  h_{\lambda_0}
  \left(
  u-x,v-y;
  D(x,y),
  \alpha_x(x,y),
  \alpha_y(x,y)
  \right).
\end{equation}

\subsection{深度分层卷积近似}

将深度离散为：
\begin{equation}
  z_1,z_2,\ldots,z_K.
\end{equation}

定义软深度 mask：
\begin{equation}
  M_k(x,y)
  =
  \frac{
  \exp\left[-\frac{(D(x,y)-z_k)^2}{2\sigma_z^2}\right]
  }{
  \sum_{j=1}^{K}
  \exp\left[-\frac{(D(x,y)-z_j)^2}{2\sigma_z^2}\right]
  }.
\end{equation}
满足：
\begin{equation}
  \sum_{k=1}^{K}M_k(x,y)=1.
\end{equation}

第 $k$ 个深度层图像：
\begin{equation}
  S_k(x,y)=M_k(x,y)I_G^{\mathrm{lin}}(x,y).
\end{equation}

若第 $k$ 层 PSF 为 $h_k(u,v)=h(u,v;z_k,\lambda_0)$，则：
\begin{equation}
  Y_k(u,v)=(h_k*S_k)(u,v).
\end{equation}

所有层相加：
\begin{equation}
\boxed{
  Y(u,v)
  =
  \sum_{k=1}^{K}
  \left[
  h_k*
  \left(
  M_k I_G^{\mathrm{lin}}
  \right)
  \right](u,v).
}
\end{equation}

展开为：
\begin{equation}
\boxed{
  Y(u,v)
  =
  \sum_{k=1}^{K}
  \sum_x\sum_y
  M_k(x,y)
  I_G^{\mathrm{lin}}(x,y)
  h_k(u-x,v-y).
}
\end{equation}

\subsection{RGB 三波长整图响应}

若采用 RGB 三代表波长近似，则：
\begin{equation}
  Y(u,v)
  =
  \sum_{c\in\{R,G,B\}}
  \sum_x\sum_y
  I_c^{\mathrm{lin}}(x,y)
  h_{\lambda_c}
  \left(
  u,v;
  x,y,
  D(x,y),
  \alpha_x(x,y),
  \alpha_y(x,y)
  \right).
\end{equation}

分层形式：
\begin{equation}
  Y(u,v)
  =
  \sum_{c\in\{R,G,B\}}
  \sum_{k=1}^{K}
  \left[
  h_{k,c}*
  \left(
  M_k I_c^{\mathrm{lin}}
  \right)
  \right](u,v),
\end{equation}
其中：
\begin{equation}
  h_{k,c}=h(u,v;z_k,\lambda_c).
\end{equation}

\subsection{CMOS 噪声与量化}

设曝光和增益合并为 $g$，理想光强 $Y(u,v)$ 对应期望电子数：
\begin{equation}
  \mu(u,v)=gY(u,v).
\end{equation}

泊松光子噪声：
\begin{equation}
  Y_{\mathrm{shot}}(u,v)\sim
  \mathrm{Poisson}(\mu(u,v)).
\end{equation}

读出噪声：
\begin{equation}
  n_r(u,v)\sim \mathcal N(0,\sigma_r^2).
\end{equation}

最终数字图像：
\begin{equation}
  \tilde Y(u,v)
  =
  Q\left[
  \mathrm{Poisson}(gY(u,v))+n_r(u,v)
  \right],
\end{equation}
其中 $Q[\cdot]$ 表示 ADC、裁剪和量化。

\section{PSF-level Fisher 信息推导}

\subsection{单参数深度估计}

对于一个点源，观测向量为：
\begin{equation}
  \bm y=
  a\bm h(z,\alpha;\phi)+\bm n,
\end{equation}
其中 $a$ 是亮度，$\bm h$ 是拉平后的 PSF，噪声为：
\begin{equation}
  \bm n\sim\mathcal N(0,\sigma^2\bm I).
\end{equation}

均值：
\begin{equation}
  \bm\mu(z)=a\bm h(z,\alpha;\phi).
\end{equation}

高斯噪声下，标量参数 $z$ 的 Fisher information 为：
\begin{equation}
  \mathcal I_z
  =
  \left(\frac{\partial\bm\mu}{\partial z}\right)^T
  \bm\Sigma^{-1}
  \left(\frac{\partial\bm\mu}{\partial z}\right).
\end{equation}

当 $\bm\Sigma=\sigma^2\bm I$：
\begin{equation}
\boxed{
  \mathcal I_z^{\mathrm{psf}}
  =
  \frac{a^2}{\sigma^2}
  \left\|
  \frac{\partial\bm h(z,\alpha;\phi)}{\partial z}
  \right\|^2 .
}
\end{equation}

对应 CRLB：
\begin{equation}
  \mathrm{Var}(\hat z)\ge
  \frac{1}{\mathcal I_z^{\mathrm{psf}}}.
\end{equation}

\subsection{多参数 nuisance-aware Fisher}

真实中深度、角度、亮度可能混淆。设：
\begin{equation}
  \bm\theta=
  \begin{bmatrix}
  z\\
  \alpha\\
  a
  \end{bmatrix}.
\end{equation}

Jacobian：
\begin{equation}
  \bm J=
  \frac{\partial\bm\mu}{\partial\bm\theta}
  =
  \begin{bmatrix}
  \frac{\partial\bm\mu}{\partial z}&
  \frac{\partial\bm\mu}{\partial\alpha}&
  \frac{\partial\bm\mu}{\partial a}
  \end{bmatrix}.
\end{equation}

其中：
\begin{equation}
  \frac{\partial\bm\mu}{\partial z}
  =
  a\frac{\partial\bm h}{\partial z},\quad
  \frac{\partial\bm\mu}{\partial\alpha}
  =
  a\frac{\partial\bm h}{\partial\alpha},\quad
  \frac{\partial\bm\mu}{\partial a}
  =
  \bm h.
\end{equation}

Fisher matrix：
\begin{equation}
  \bm F=\bm J^T\bm\Sigma^{-1}\bm J.
\end{equation}

若 $\bm\theta=[z,\bm\eta^T]^T$，其中 $\bm\eta$ 是干扰参数，则分块：
\begin{equation}
  \bm F=
  \begin{bmatrix}
  F_{zz}&\bm F_{z\eta}\\
  \bm F_{\eta z}&\bm F_{\eta\eta}
  \end{bmatrix}.
\end{equation}

只关心深度时，有效 Fisher 为 Schur complement：
\begin{equation}
\boxed{
  \mathcal I_z^{\mathrm{eff}}
  =
  F_{zz}
  -
  \bm F_{z\eta}
  \bm F_{\eta\eta}^{-1}
  \bm F_{\eta z}.
}
\end{equation}

该式表示：从总深度敏感性中扣除与角度、亮度、波长等干扰参数混淆的部分。

\subsection{泊松噪声下的 Fisher}

若 CMOS 光子计数近似：
\begin{equation}
  y_i\sim\mathrm{Poisson}(\mu_i(\bm\theta)),
\end{equation}
则：
\begin{equation}
\boxed{
  F_{jk}
  =
  \sum_i
  \frac{1}{\mu_i}
  \frac{\partial\mu_i}{\partial\theta_j}
  \frac{\partial\mu_i}{\partial\theta_k}.
}
\end{equation}

只估计深度：
\begin{equation}
\boxed{
  \mathcal I_z
  =
  \sum_i
  \frac{1}{\mu_i}
  \left(
  \frac{\partial\mu_i}{\partial z}
  \right)^2 .
}
\end{equation}

这说明 PSF 变化大有用，但光强过低会因 photon noise 降低有效信息。

\section{Image-level Fisher 信息推导}

\subsection{整图观测模型}

将 CMOS 图像拉平成向量：
\begin{equation}
  \bm y\in\mathbb R^N.
\end{equation}
将深度图拉平成：
\begin{equation}
  \bm D=[D_1,D_2,\ldots,D_M]^T.
\end{equation}
给定 RGB 图 $\bm X$，超表面编码图均值为：
\begin{equation}
  \bm\mu(\bm D,\bm X;\phi)=\mathcal H_\phi(\bm X,\bm D).
\end{equation}
观测模型：
\begin{equation}
  \bm y=
  \bm\mu(\bm D,\bm X;\phi)+\bm n,\qquad
  \bm n\sim\mathcal N(0,\sigma^2\bm I).
\end{equation}

\subsection{均值函数的离散形式}

令 $j$ 表示输入场景像素，$i$ 表示 CMOS 像素。则：
\begin{equation}
  \mu_i(\bm D,\bm X;\phi)
  =
  \sum_{j=1}^{M}
  X_j
  h_i(j,D_j;\phi).
\end{equation}

这里 $h_i(j,D_j;\phi)$ 表示第 $j$ 个场景点对第 $i$ 个 CMOS 像素的 PSF 贡献。

\subsection{深度图 Fisher matrix}

Jacobian：
\begin{equation}
  \bm J_D=
  \frac{\partial\bm\mu}{\partial\bm D},
  \qquad
  [\bm J_D]_{ij}
  =
  \frac{\partial\mu_i}{\partial D_j}.
\end{equation}

由：
\begin{equation}
  \mu_i=
  \sum_{j=1}^{M}
  X_j h_i(j,D_j;\phi)
\end{equation}
得：
\begin{equation}
\boxed{
  [\bm J_D]_{ij}
  =
  X_j
  \frac{\partial h_i(j,D_j;\phi)}{\partial D_j}.
}
\end{equation}

高斯噪声下整图 Fisher：
\begin{equation}
\boxed{
  \bm F_D^{\mathrm{img}}
  =
  \bm J_D^T
  \bm\Sigma^{-1}
  \bm J_D.
}
\end{equation}

若 $\bm\Sigma=\sigma^2\bm I$：
\begin{equation}
\boxed{
  \bm F_D^{\mathrm{img}}
  =
  \frac{1}{\sigma^2}
  \bm J_D^T\bm J_D.
}
\end{equation}

元素展开：
\begin{equation}
  F_{pq}
  =
  \frac{1}{\sigma^2}
  \sum_i
  \frac{\partial\mu_i}{\partial D_p}
  \frac{\partial\mu_i}{\partial D_q}.
\end{equation}

代入：
\begin{equation}
\boxed{
  F_{pq}
  =
  \frac{X_pX_q}{\sigma^2}
  \sum_i
  \frac{\partial h_i(p,D_p;\phi)}{\partial D_p}
  \frac{\partial h_i(q,D_q;\phi)}{\partial D_q}.
}
\end{equation}

\subsection{为什么 image-level Fisher 更真实}

PSF-level Fisher 只关心点源：
\begin{equation}
  \frac{\partial\bm h}{\partial z}.
\end{equation}
Image-level Fisher 关心：
\begin{equation}
  \frac{\partial\bm\mu}{\partial D_j},
\end{equation}
即一个场景像素的深度变化会怎样改变整张 CMOS 图像。它包含亮度、纹理、PSF 重叠、低纹理区域、遮挡边界和自然图像叠加效应，因此更接近最终任务。

\subsection{实际可计算近似}

完整 $\bm F_D$ 维度巨大，实际可用以下近似。

\paragraph{对角 Fisher}
只看每个像素深度自身的信息：
\begin{equation}
  F_{jj}
  =
  \frac{1}{\sigma^2}
  \left\|
  \frac{\partial\bm\mu}{\partial D_j}
  \right\|^2.
\end{equation}
代入：
\begin{equation}
  F_{jj}
  =
  \frac{X_j^2}{\sigma^2}
  \left\|
  \frac{\partial\bm h(j,D_j;\phi)}{\partial D_j}
  \right\|^2.
\end{equation}

\paragraph{Patch-level Fisher}
在局部 patch 内，用低维参数描述深度，例如平面：
\begin{equation}
  D_j(a,b,c)=a x_j+b y_j+c.
\end{equation}
参数为 $\bm\theta=[a,b,c]^T$。局部 Fisher：
\begin{equation}
  \bm F_\Omega=
  \bm J_\Omega^T\bm\Sigma^{-1}\bm J_\Omega.
\end{equation}

\paragraph{Depth-bin Fisher}
将图像分成深度层：
\begin{equation}
  \bm\mu=
  \sum_{k=1}^{K}
  \bm H_k(\phi)\bm s_k,
\end{equation}
其中 $\bm s_k=M_k(\bm D)\odot\bm X$。可以分析不同深度层响应之间的相关性，避免不同深度层编码过于相似。

\section{优化目标整理}

\subsection{PSF-level 信息优化}

单纯最大化深度敏感性：
\begin{equation}
  \phi^\star
  =
  \arg\max_\phi
  \mathbb E_{z,\alpha}
  \left[
  \log
  \mathcal I_z^{\mathrm{psf}}(z,\alpha;\phi)
  \right].
\end{equation}

考虑干扰参数：
\begin{equation}
\boxed{
  \phi^\star
  =
  \arg\max_\phi
  \mathbb E_{z,\alpha}
  \left[
  \log
  \mathcal I_z^{\mathrm{eff}}(z,\alpha;\phi)
  \right].
}
\end{equation}

强调长距离和边缘视场最坏情况：
\begin{equation}
\boxed{
  \phi^\star
  =
  \arg\max_\phi
  \min_{z\in\mathcal Z,\alpha\in\mathcal A}
  \mathcal I_z^{\mathrm{eff}}(z,\alpha;\phi).
}
\end{equation}

\subsection{Image-level Fisher 优化}

对整图或 patch 近似：
\begin{equation}
  \mathcal L_{\mathrm{imgFI}}
  =
  -
  \mathbb E_{(\bm X,\bm D)}
  \left[
  \sum_j
  w_j
  \log
  \left(
  \left\|
  \frac{\partial\bm\mu_\phi}{\partial D_j}
  \right\|^2
  +
  \epsilon
  \right)
  \right].
\end{equation}

联合 PSF 与图像级目标：
\begin{equation}
  \phi^\star
  =
  \arg\min_\phi
  \left[
  \lambda_{\mathrm{psf}}\mathcal L_{\mathrm{psf}}
  +
  \lambda_{\mathrm{img}}\mathcal L_{\mathrm{imgFI}}
  +
  \lambda_{\mathrm{fab}}\mathcal L_{\mathrm{fab}}
  \right].
\end{equation}

\subsection{恢复网络训练}

优化得到相位 $\phi^\star$ 后，生成超表面编码图：
\begin{equation}
  \bm Y_{\phi^\star}
  =
  \mathcal H_{\phi^\star}(\bm X,\bm D).
\end{equation}

训练深度恢复网络：
\begin{equation}
  \hat{\bm D}
  =
  f_\psi(\bm X,\bm Y_{\phi^\star},\bm R_{\mathrm{ray}}),
\end{equation}
其中 $\bm R_{\mathrm{ray}}$ 是每个像素的光线方向或视场角编码。

监督损失：
\begin{equation}
  \mathcal L_{\mathrm{depth}}
  =
  \|M\odot(\hat{\bm D}-\bm D)\|_1.
\end{equation}

可加入梯度损失：
\begin{equation}
  \mathcal L_{\mathrm{grad}}
  =
  \|\nabla\hat{\bm D}-\nabla\bm D\|_1.
\end{equation}

若使用 Depth Anything V2 等 teacher，可加入蒸馏：
\begin{equation}
  \mathcal L_{\mathrm{distill}}
  =
  \left\|
  \mathrm{SSI}(\hat{\bm D})
  -
  \mathrm{SSI}(\bm D_T)
  \right\|_1.
\end{equation}

最终训练：
\begin{equation}
  \psi^\star
  =
  \arg\min_\psi
  \left[
  \lambda_1\mathcal L_{\mathrm{depth}}
  +
  \lambda_2\mathcal L_{\mathrm{grad}}
  +
  \lambda_3\mathcal L_{\mathrm{distill}}
  \right].
\end{equation}

\section{推荐实施顺序}

\begin{enumerate}
  \item 固定单波长 $\lambda_0$、孔径 $D_{\mathrm{ap}}$、超表面-CMOS 距离 $d$、像素尺寸和目标深度范围。
  \item 实现点源 PSF 计算：
  \begin{equation}
    h=
    \left|
    \mathcal P_d
    \left[
    P t_\phi U_{\mathrm{in}}
    \right]
    \right|^2.
  \end{equation}
  \item 计算 $\mathcal I_z^{\mathrm{psf}}$ 和 $\mathcal I_z^{\mathrm{eff}}$，画出 CRLB 随距离和视场角变化图。
  \item 用低维相位参数化：
  \begin{equation}
    \phi(x,y)=\sum_m c_m B_m(x,y),
  \end{equation}
  优化系数 $\bm c$。
  \item 扫描口径 $D_{\mathrm{ap}}$ 和传播距离 $d$，画 feasibility map。
  \item 用 RGB-D 数据和分层 PSF 公式生成超表面编码图：
  \begin{equation}
    Y=
    \sum_k
    h_k*(M_kI_G^{\mathrm{lin}}).
  \end{equation}
  \item 训练小型 decoder 或基于 Depth Anything V2 的 prompt/fine-tuning 网络。
  \item 做严格消融：普通透镜、随机相位、人工旋转/DH PSF、Fisher 相位、Fisher+image-level 相位。
  \item 后续再引入相位量化、透过率、RCWA meta-atom library、偏振和宽谱响应。
\end{enumerate}

\section{一句话总结}

本项目的公式主线是：
\begin{equation}
\begin{aligned}
  \text{像素}&\xrightarrow{K^{-1}}\text{光线方向},\\
  \text{RGB}&\xrightarrow{\mathrm{linearize}}\text{入射强度},\\
  \text{点源}&\xrightarrow{P t_\phi+\mathcal P_d}\text{PSF},\\
  \text{整图}&\xrightarrow{\sum \text{PSF splatting}}\text{CMOS 图},\\
  \text{Fisher}&\xrightarrow{\max}\text{信息最优相位}.
\end{aligned}
\end{equation}

最值得坚持的核心不是简单 ``设计一个特殊 PSF''，而是：
\begin{quote}
用 Fisher/CRLB 在物理测量层面优化 CMOS 图像对深度的有效可观测性，再用深度学习模型从这个更有信息的物理编码图中恢复 metric depth。
\end{quote}

\end{document}
